% (The MIT License)
%
% Copyright (c) 2023-2024 Yegor Bugayenko
%
% Permission is hereby granted, free of charge, to any person obtaining a copy
% of this software and associated documentation files (the 'Software'), to deal
% in the Software without restriction, including without limitation the rights
% to use, copy, modify, merge, publish, distribute, sublicense, and/or sell
% copies of the Software, and to permit persons to whom the Software is
% furnished to do so, subject to the following conditions:
%
% The above copyright notice and this permission notice shall be included in all
% copies or substantial portions of the Software.
%
% THE SOFTWARE IS PROVIDED 'AS IS', WITHOUT WARRANTY OF ANY KIND, EXPRESS OR
% IMPLIED, INCLUDING BUT NOT LIMITED TO THE WARRANTIES OF MERCHANTABILITY,
% FITNESS FOR A PARTICULAR PURPOSE AND NONINFRINGEMENT. IN NO EVENT SHALL THE
% AUTHORS OR COPYRIGHT HOLDERS BE LIABLE FOR ANY CLAIM, DAMAGES OR OTHER
% LIABILITY, WHETHER IN AN ACTION OF CONTRACT, TORT OR OTHERWISE, ARISING FROM,
% OUT OF OR IN CONNECTION WITH THE SOFTWARE OR THE USE OR OTHER DEALINGS IN THE
% SOFTWARE.

\documentclass{article}
\usepackage{../lecture-notes/notes}
\newcommand*\thetitle{Reflection}
\newcommand*\thesubtitle{Classpath, Casting, Annotations}
\begin{document}

\plush{\lnTitlePage{7}{8}{vYw_GcDkPQA}}

\pptToc

\plush{\pptChapter[What?]{What Is Reflection?}}

\pptSection[Example]{Simple Example of Reflection in Java}
\begin{multicols}{2}
{\small\begin{ffcode}
void print(Object x) {
  if (x instance of File) {
    System.out.println("Yes!");
  }
}
\end{ffcode}
}
\par\columnbreak\par
``Reflective computational systems allow computations to observe and
modify properties of their own behavior. It would be desirable for computations to avail themselves of these
reflective capabilities, examining themselves in order to make use of
\ul{meta-level information} in decisions about what to do next.''
\lnSource{sobel1996introduction}
\end{multicols}
\par
\plush{}

\lnQuote
  [Yue Li]
  {yue-li}
  {As reflection is increasingly used in Java programs, the cost of imprecise reflection handling has increased dramatically... Almost all the analyses reported are \ul{unsound} in the presence of reflection since it is either \ul{ignored} or handled \ul{partially}.}
  {li2014self}

\plush{\pptChapter[Casting]{Type Casting and Subsumption}}

\pptSection[Iterable]{Iterable \(\to\) Collection}
\begin{pptWide}{2}
Downcasting (\textcolor{red}{wrong!}):\par
{\small\begin{ffcode}
int sizeOf(Iterable items) {
  int size = 0;
  if (items (*@\textcolor{orange}{instanceof}@*) Collection) {
    size = (*@\textcolor{orange}{((Collection)}@*) items).size();
  } else {
    for (Object item : items) {
      ++size;
    }
  }
  return size;
}
\end{ffcode}
}
\par\columnbreak\par
Method overloading (\textcolor{green}{right!}):\par
{\small\begin{ffcode}
int sizeOf(Iterable items) {
  int size = 0;
  for (Object item : items) {
    ++size;
  }
  return size;
}

int sizeOf((*@\textcolor{orange}{Collection}@*) items) {
  return items.size();
}
\end{ffcode}
}
\end{pptWide}
\par
\plush{}

\pptSection[Coupling]{Implicit Coupling}
\begin{pptWide}{2}
{\small\begin{ffcode}
int sizeOf(Iterable items) {
  int size = 0;
  (*@\textcolor{orange}{if}@*) (items instanceof Collection) {
    size = ((Collection) items).size();
  } (*@\textcolor{orange}{else}@*) {
    for (Object item : items) {
      ++size;
    }
  }
  return size;
}
\end{ffcode}
}
\par\columnbreak\par
{\small\begin{ffcode}
Iterable<Book> books1 =
  new PgBooks1("localhost:5432");
int s1 = (*@\textcolor{orange}{sizeOf}@*)(books1);

Collection<Book> books2 =
  new PgBooks2("localhost:5432");
int s2 = (*@\textcolor{orange}{sizeOf}@*)(books2);
\end{ffcode}
}
\end{pptWide}
\par
It may be hard to understand why |s2| is evaluated much faster than |s1|,
while the signature of |sizeOf()| is the same~\citep{bugayenko2015blog0402}.
\plush{}

\pptSection[JEP-394]{Pattern Matching in Java 16}%
\begin{pptWide}{2}%
Java 11 (\textcolor{red}{wrong!}):
\par
{\small\begin{ffcode}
int sizeOf(Iterable items) {
  int size = 0;
  if (items (*@\textcolor{orange}{instanceof}@*) Collection) {
    size = ((*@\textcolor{orange}{(Collection)}@*) items).size();
  } else {
    for (Object item : items) {
      ++size;
    }
  }
  return size;
}
\end{ffcode}
}
\par\columnbreak\par
\pptPin{\raggedleft\includegraphics[width=.3\linewidth]{java-logo.pdf}}%
Java 16 (\textcolor{red}{\textbf{even worse!}}):
\par
{\small\begin{ffcode}
int sizeOf(Iterable items) {
  int size = 0;
  if ((*@\textcolor{orange}{items instanceof Collection c}@*)) {
    size = c.size();
  } else {
    for (Object item : items) {
      ++size;
    }
  }
  return size;
}
\end{ffcode}
}
\end{pptWide}
\par
\plush{}

\pptSection[Matching]{C\#, Rust, and pattern matching}
\begin{pptWide}{2}
C\#:\par
{\scriptsize\begin{ffcode}
public int sizeOf<T>(IEnumerable<T> items) {
  if (items (*@\textcolor{orange}{is}@*) IList<T> list) {
    return list.Count;
  } else {
    return // count them one by one
  }
}
\end{ffcode}
}
\par
Some other languages have \ul{pattern matching} feature,
including Kotlin, Scala, Haskell, Elixir, Swift, F\#, and Erlang, ...
which \textcolor{red}{contradicts} the principle of encapsulation.
\par\columnbreak\par
Rust:
\par
{\scriptsize\begin{ffcode}
enum Color {
  RGB(u8, u8, u8),
  Transparent
}
fn paint(c: Color) {
  (*@\textcolor{orange}{match}@*) c {
    Color::RGB(r, g, b) =>
      println!("#{r}{g}{b}"),
    Color::Transparent =>
      println!("none")
  }
}
fn main() {
  let c = Color::RGB(64, 16, 0);
  paint(c);
}
\end{ffcode}
}
\end{pptWide}
\par
\plush{}

\plush{\pptChapter[Factory]{Factory Method}}

\pptSection[IF]{Conditional object construction}
\begin{pptWide}{2}
This is \textcolor{red}{wrong}:
\par
{\scriptsize\begin{ffcode}
interface Figure
  int surface();
class Square implements Figure
class Triangle implements Figure
class Polygon implements Figure

class FactoryOfFigures
  Figure make(int sides) {
    if (sides == 3) {
      return (*@\textcolor{orange}{new Triangle()}@*);
    } else if (sides == 4) {
      return (*@\textcolor{orange}{new Square()}@*);
    } else {
      return (*@\textcolor{orange}{new Polygon(sides)}@*);
    }
  }
\end{ffcode}
}
\par\columnbreak\par
This is \textcolor{green}{better}:
\par
{\scriptsize\begin{ffcode}
class PolymorphicFigure implements Figure
  PolymorphicFigure(int sides)
  @Override int surface() {
    if (sides == 3) {
      return (*@\textcolor{orange}{new Triangle().surface()}@*);
    } else if (sides == 4) {
      return (*@\textcolor{orange}{new Square().surface()}@*);
    } else {
      return (*@\textcolor{orange}{new Polygon(sides).surface()}@*);
    }
  }
\end{ffcode}
}
\par
Here, the semantic of object construction is not visible to
the client---the \ul{coupling} is ``loose''~\citep{bugayenko2022blog0605}.
\end{pptWide}
\plush{}

\pptSection[forName]{Generating class name from a string}
\begin{pptWide}{2}
This is \textcolor{red}{wrong}:
\par
{\scriptsize\begin{ffcode}
interface Figure
  int surface();

class Square implements Figure
class Triangle implements Figure
class Polygon implements Figure

class FactoryOfFigures
  Figure make(String name) throws Exception {
    Class<?> c = (*@\textcolor{orange}{Class.forName}@*)(name);
    return c.getConstructor()(*@\textcolor{orange}{.newInstance()}@*);
  }
\end{ffcode}
}
\par\columnbreak\par
This is \textcolor{green}{better}:
\par
{\scriptsize\begin{ffcode}
class PolymorphicFigure implements Figure
  PolymorphicFigure(String name)
  @Override int surface() {
    if ((*@\textcolor{orange}{name.equals}@*)("Triangle")) {
      return new Triangle().surface();
    } else if ((*@\textcolor{orange}{name.equals}@*)("Square")) {
      return new Square().surface();
    } else {
      return new Polygon(?).surface();
    }
  }
\end{ffcode}
}
\par
This is better since the mechanics of class finding
is explicit---no surprises expected~\citep{bugayenko2017blog1114}.
\end{pptWide}
\plush{}

\plush{\pptChapter[Classpath]{Classpath Scanning}}

\pptSection[Class]{Finding Java classes}
\begin{pptWide}{2}
{\small\begin{ffcode}
interface Foo {}

class Bar implements Foo {}

Reflections rts =
  (*@\textcolor{orange}{new Reflections}@*)("");
Set<Class<?>> types = rts.get(
  SubTypes.of(Foo.class).asClass()
);
\end{ffcode}
}
\par\columnbreak\par
{\small\begin{ffcode}
public @interface Foo {}

@Foo
class Bar {}

Reflections rts =
  new Reflections("");
Set<Class<?>> types = rts.get(
  SubTypes.of(
    (*@\textcolor{orange}{TypesAnnotated.with}@*)(Foo.class)
  ).asClass()
);
\end{ffcode}
}
\end{pptWide}
\par
The library is called \href{https://github.com/ronmamo/reflections}{Reflections}.
Instead, use explicit object instantiation.
\plush{}

\plush{\pptChapter{Annotations}}

\pptSection[Class]{I lieu of static methods}
\begin{pptWide}{2}
{\scriptsize\begin{ffcode}
interface Pub
  String isbn();

class Book implements Pub
  @Override public String isbn()
    /* ... */
  public (*@\textcolor{orange}{static}@*) String category()
    return "book";

class Journal implements Pub
  @Override public String isbn()
    /* ... */
  public (*@\textcolor{orange}{static}@*) String category()
    return "journal";
\end{ffcode}
}
\par\columnbreak\par
{\scriptsize\begin{ffcode}
interface Pub
  String isbn();

@Target(ElementType.CLASS)
@Retention(RetentionPolicy.SOURCE)
public @interface Category
  String value();

(*@\textcolor{orange}{@Category("book")}@*)
class Book implements Pub
  @Override public String isbn()
    /* ... */

(*@\textcolor{orange}{@Category("journal")}@*)
class Journal implements Pub
  @Override public String isbn()
    /* ... */
\end{ffcode}
}
\end{pptWide}
\par
\plush{}

\pptSection[Method]{Locating methods}
\begin{pptWide}{2}
{\small\begin{ffcode}
@Target(ElementType.METHOD)
@Retention(RetentionPolicy.SOURCE)
public @interface Path
  String url;

class BookController
  (*@\textcolor{orange}{@Path("/book-title")}@*)
  String title()
    // Build HTML page and return it
\end{ffcode}
}
\par\columnbreak\par
{\scriptsize\begin{ffcode}
String dispatch(String url) {
  c = new BooksController();
  b = BooksController.class;
  for (Method m : b(*@\textcolor{orange}{.getDeclaredMethods()}@*)) {
    if (m.isAnnotationPresent(Path.class)) {
      Annotation a = m(*@\textcolor{orange}{.getAnnotation}@*)(Path.class);
      if (a(*@\textcolor{orange}{.url()}@*).equals(url)) {
        return m.invoke(c);
      }
    }
  }
  return "404 Page not found";
}
\end{ffcode}
}
\end{pptWide}
\par
The |BookContoller| doesn't know how exactly the data
in its annotation is being used and by whom~\citep{bugayenko2016blog0412}.
\plush{}

\pptSection[DIC]{Dependency Injection Container}
\begin{pptWide}{2}
{\scriptsize\begin{ffcode}
interface Shipment
  int cost();

class Cart
  (*@\textcolor{orange}{@Inject}@*) private Shipment shmt;
  private Book book;
  void setBook(Book b)
    this.book = b;
  int cost()
    return this.book.price() + this.shmt.cost();

container = new Container();
c = container(*@\textcolor{orange}{.make(Cart.class)}@*);
c.setBook(new Book("1984"));
x = c.cost();
\end{ffcode}
}
\par\columnbreak\par
{\scriptsize\begin{ffcode}
class Container {
  private HashMap<Class, Object> cache =
    new ConcurrentHashMap<>();
  T make(Class<T> type) {
    // 1. Find @Inject-annotated "shmt" field;
    // (*@\textcolor{orange}{2. Make an instance of "Shipment";}@*)
    // 3. Store it in the "cache";
    // 4. Make an instance of "Cart";
    // 5. Store "cart" in the "cache";
    // 6. Assign "shipment" to "cart.shmt";
    // 7. Return "cart".
  }
}
\end{ffcode}
}
How do you think, at the step no.2,
what class will be instantiated?~\citep{bugayenko2014blog1003}
\end{pptWide}
\par
\plush{}

\pptSection[DI]{Dependency Injection \textit{without} a Container}
\begin{pptWide}{2}
{\scriptsize\begin{ffcode}
interface Shipment
  int cost();

class Cart
  (*@\textcolor{orange}{@Inject}@*) private Shipment shmt;
  private Book book;
  void setBook(Book b)
    this.book = b;
  int cost()
    return this.book.price() + this.shmt.cost();

container = new Container();
c = container.make(Cart.class);
c.setBook(new Book("1984"));
x = c.cost();
\end{ffcode}
}
\par\columnbreak\par
{\scriptsize\begin{ffcode}
interface Shipment
  int cost();

class Cart
  private (*@\textcolor{orange}{final}@*) Shipment shmt;
  private (*@\textcolor{orange}{final}@*) Book book;
  Cart(Shipment s, Book b)
    this.shmt = s;
    this.book = b;
  int cost()
    return this.book.price() + this.shmt.cost();

c = (*@\textcolor{orange}{new Cart(new MyShipment(), new Book("1984"))}@*);
x = c.cost();
\end{ffcode}
}
\end{pptWide}
\par
\plush{}

\plush{\pptChapter[Types]{Discrimination by Type}}

\pptSection[Polymorphism]{Polymorphism vs. Casting}
\begin{pptWide}{2}
{\small\begin{ffcode}
interface Figure
  void rotate(int d);
class Circle implements Figure
  void rotate(int d) //...
  int radius() //...
class Square implements Figure
  void rotate(int d) //...
  int side() //...
\end{ffcode}
}
\par
Philosophically speaking, type casting is discrimination~\citep{bugayenko2020blog1110}.
\par\columnbreak\par
{\small\begin{ffcode}
// This is polymorphism:
int surface(Figure f)
  return f(*@\textcolor{orange}{.surface()}@*);

// This is type casting:
int surface(Figure f)
  if (f instanceof Circle c) {
    return c(*@\textcolor{orange}{.radius()}@*)
  } else if (f instanceof Square s) {
    return s(*@\textcolor{orange}{.side()}@*) * s.side();
  } else {
    throw new Exception("oops");
  }
\end{ffcode}
}
\end{pptWide}
\par
\plush{}

\end{document}
